% تمرین نمونه است تا مخزن از روز نخست خالی نباشد؛ پیش از استفاده بازبینی‌اش کنید.
\documentclass{../is-assignment}

\عنوان{\متن‌لاتین{Web Hardening}}
\ترم{پاییز ۱۴۰۵}

\begin{document}

\عنوان‌ساز

\فهرست‌مطالب

\قسمت{مقدمه}

در این تمرین یک برنامه‌ی وب کوچک و آسیب‌پذیر را می‌گیرید، ضعف‌هایش را پیدا می‌کنید و می‌بندید.
برنامه می‌تواند یکی از نمونه‌های آموزشی شناخته‌شده باشد یا برنامه‌ی ساده‌ای که خودتان می‌نویسید،
به شرطی که دست‌کم ورود کاربر، نشست و یک فرم داشته باشد.

\شروع{نکته}
همه‌ی کار روی سامانه‌ی خودتان انجام می‌شود.
آزمودن این حمله‌ها روی سامانه‌ای که اجازه‌اش را ندارید، تخلف است و نمره‌ی این تمرین را صفر می‌کند.
\پایان{نکته}

\قسمت{نشست}

\شروع{شمارش}
\فقره کوکی نشست برنامه را ببینید و بگویید کدام صفت‌ها را ندارد.
\فقره \متن‌لاتین{HttpOnly}، \متن‌لاتین{Secure} و \متن‌لاتین{SameSite} را اضافه کنید و برای هرکدام بگویید جلوی چه چیزی را گرفتید.
\فقره نشان دهید پس از خروج کاربر، شناسه‌ی نشست دیگر پذیرفته نمی‌شود.
\پایان{شمارش}

\نمره{۳}

\قسمت{تزریق و بازتاب}

\شروع{شمارش}
\فقره یک \متن‌لاتین{XSS} بازتابی روی برنامه‌ی خودتان نشان دهید و سپس آن را ببندید.
      توضیح دهید چرا پالایش ورودی به تنهایی کافی نیست و \متن‌لاتین{escape} در خروجی لازم است.
\فقره اگر برنامه پرس‌وجوی \متن‌لاتین{SQL} دارد، نشان دهید چگونه با \متن‌لاتین{prepared statement} از تزریق جلوگیری می‌کنید.
\پایان{شمارش}

\نمره{۴}

\قسمت{ارتباط امن}

\شروع{شمارش}
\فقره برنامه را پشت \متن‌لاتین{HTTPS} بیاورید (گواهی خودامضا برای این تمرین بس است).
\فقره سرآیند \متن‌لاتین{HSTS} را اضافه کنید و بگویید چه چیزی را حل می‌کند که تنها روشن کردن \متن‌لاتین{HTTPS} حل نمی‌کند.
\پایان{شمارش}

\نمره{۳}

\شروع{امتیازی}
یک سیاست \متن‌لاتین{Content-Security-Policy} بنویسید که \متن‌لاتین{XSS} بخش پیشین را حتی بدون اصلاح کد بی‌اثر کند،
و بگویید چرا \متن‌لاتین{unsafe-inline} آن را بی‌معنا می‌کند.
\پایان{امتیازی}

\قسمت{تحویل}

مخزن کد به همراه گزارشی که برای هر ضعف سه چیز را نشان دهد: پیش از اصلاح، تغییری که دادید، و پس از اصلاح.

\قوانین

\پایان‌ساز

\end{document}
